\section{Related Work (extended)}
\label{sec:related}

\subsection{Fragmentomics}
The cfDNA fragmentomics literature is the empirical foundation for HAVI-Methyl's input feature set. \citet{snyder2016} mapped nucleosome positioning from cfDNA and established that fragment-length distributions encode tissue-of-origin information through nucleosome occupancy patterns; the canonical structure is a dominant peak at approximately $167$ bp (the chromatosome) with $\sim 10.4$ bp periodicity in the $100$--$160$ bp range, plus weaker di- and tri-nucleosome modes. \citet{sun2019orientationaware} introduced orientation-aware fragmentation and demonstrated that the relative orientation of paired-end reads near transcription start sites is informative for tissue identification. \citet{zhou2022} characterized end-motif distributions in healthy and cancer plasma and showed that the 4-mer at fragment cut sites distinguishes cancer from healthy at the population level. \citet{cristiano2019delfi} introduced DELFI as the first genome-wide fragmentation profile classifier for cancer detection, and \citet{mathios2021delfi} extended DELFI to lung-cancer detection at AUC$\approx 0.94$ on $n{=}365$ samples. \citet{esfahani2022epicseq} demonstrated that fragmentation entropy at gene promoters predicts gene expression, establishing that the fragmentomic signal carries information well beyond methylation alone. \citet{zviran2020mrdetect} and \citet{widman2024mrdedge} addressed minimal-residual-disease monitoring with tumor-informed mutational integration and machine-learning-guided signal enrichment, providing a contrast point for our tumor-naive methylation-prediction route. Crucially for our framing, \citet{han2024methfrag} established directly that DNA methylation modulates cfDNA fragmentation: methylation status shifts the cut-site preferences and length distribution at a measurable level, providing the mechanistic justification for predicting methylation from fragmentomics. FinaleMe \citep{liu2024finaleme} synthesized these signals into a per-sample HMM that recovers methylation. HAVI-Methyl takes the same input feature set but treats fragments as observations of a hierarchical latent-variable model rather than emissions of a per-sample HMM, which is the central methodological difference.

\subsection{Variational Bayes}
\citet{mackay2003} and \citet{beal2003} laid the foundations of variational Bayes for graphical models. \citet{hoffman2013svi} introduced stochastic VI for conditionally conjugate models and is the basis for our population-layer update. \citet{hoffmanblei2015ssvi} extended SVI to structured variational families. \citet{titsias2014doublystoch} and \citet{kucukelbir2017advi} introduced doubly stochastic VI and ADVI for non-conjugate models. \citet{kingma2014vae,rezende2014} introduced amortized VAEs and reparameterization gradients. \citet{rezende2015flows,kingma2016iaf,durkan2019nsf} developed normalizing flows of progressively increasing expressiveness, and recent work has extended this in two directions relevant to the methylation $\beta\in[0,1]$ support: \citet{mathieu2020rcnf} introduced Riemannian continuous normalizing flows on bounded manifolds, and \citet{wu2020snf} introduced stochastic normalizing flows that interleave invertible and stochastic transformations. \citet{kingma2021vdm} introduced variational diffusion models with an SNR-parameterized variational lower bound that is theoretically tight. \citet{burda2016iwae} introduced IWAE; \citet{rainforth2018iwaeSNR} identified the IWAE encoder gradient SNR problem; and \citet{tucker2019dreg} resolved it with the doubly-reparameterized gradient estimator. \citet{khan2017cvi,lin2019natgrad} extended natural-gradient SVI to non-conjugate models and provided the theoretical basis for our Proposition~\ref{prop:svi}, with \citet{khan2018vadam} providing the practical Adam-based implementation route that we follow. \citet{alemi2017vib} introduced the variational information bottleneck. \citet{cremer2018amortization} introduced the amortization-gap concept that we monitor at evaluation time. \citet{blei2017vi} reviewed VI for statisticians.

\subsection{Causal and instrumental epigenetics}
\citet{ying2024causality} used Mendelian randomization on aging clocks and motivated the use of mQTLs as instrumental anchors, and \citet{ying2024causage} extended this with the CausAge / DamAge / AdaptAge family of causality-enriched clocks that decompose age-associated methylation into damage and adaptation components, motivating evaluation of HAVI-Methyl on these CpG panels specifically. \citet{min2021mqtl} catalogued mQTLs in the GoDMC consortium, providing approximately $248{,}607$ independent cis-mQTLs at $5\times 10^{-8}$ genome-wide significance and the external IV summary statistics used in our anchor loss. \citet{sanderson2022mr} reviewed Mendelian randomization methodology in a clinical-translation context. \citet{newey2003instrumental} and \citet{darolles2011nonparametric} provide the nonparametric IV completeness theory underlying Proposition~\ref{prop:identify-restated}. For counterfactual augmentation, \citet{pawlowski2020dscm} provides the deep-structural-causal-model template that we adapt for cfDNA, and \citet{kaddour2022causalml} surveys the broader causal-machine-learning landscape; \citet{lotfollahi2023cpa} provides a single-cell counterfactual-perturbation approach that informs our matched-pair construction. \citet{higginschen2022pcclocks} introduced PC-clocks and the Shrout-Fleiss ICC reliability convention adopted in our benchmarking metrics, and \citet{horvath2013,levine2018phenoage,lu2019grimage,lu2022grimage2,belsky2022dunedinpace} are the foundational and contemporary epigenetic-clock references whose downstream pipelines consume HAVI-Methyl's $\beta$ output.

\subsection{Methylation foundation models and long-range DNA encoders}
\citet{cpgpt2024,methylgpt2024} pretrained large transformers on methylation array data and produce powerful CpG embeddings, but require array-measured $\beta$ as input rather than producing it from raw WGS. Long-range DNA encoders are converging on subquadratic-attention architectures: \citet{nguyen2023hyenadna} demonstrated context up to $1$M tokens at single-nucleotide resolution; \citet{schiff2024caduceus} provided a bi-directional reverse-complement-equivariant Mamba variant; \citet{nguyen2024evo} scaled this further to a $7$B-parameter StripedHyena trained at $131{,}072$-token context; \citet{dallatorre2025nt} provides a $500$M-parameter transformer benchmarked on human-genomics tasks; \citet{zhou2024dnabert2} introduced DNABERT-2 with byte-pair-encoded tokenization and ALiBi position bias; and \citet{benegas2023gpn} showed that a convolutional masked-language-modelling encoder predicts genome-wide variant effects. These provide drop-in alternatives for the frozen sequence-context embedding in our encoder, with the empirical question of which long-range encoder gives the best methylation-prediction performance left for the benchmark plan in Section~\ref{sec:benchmark}.

\subsection{Probabilistic methylation prediction from sequencing}
The line of work most directly comparable to HAVI-Methyl is probabilistic methylation prediction from sequencing reads. DeepCpG \citep{angermueller2017deepcpg} uses a CNN over a $1001$-bp sequence window combined with a bidirectional GRU over neighbouring CpG states to predict single-cell methylation, but produces point estimates without principled uncertainty quantification or hierarchical pooling across samples. METHimpute \citep{taudt2018methimpute} uses a two-state HMM to impute missing methylation but operates on bisulfite data rather than fragmentomics, and produces a per-CpG mean rather than a calibrated posterior. MethylBERT \citep{lee2018methylbert} uses a transformer to classify methylation read-by-read and supports tumour deconvolution, but does not produce locus-level posteriors. CpGenie \citep{zeng2017cpgenie} predicts the impact of non-coding variants on methylation through a sequence-only CNN. Closer in spirit to our hierarchical Bayesian framing, BPRMeth \citep{kapourani2018bprmeth} models methylation profiles with a Bernoulli-probit regression and variational Bayes, and Melissa \citep{kapourani2019melissa} extends this with a Bayesian hierarchical model for clustering and imputation across single-cell methylomes. Melissa is the closest philosophical antecedent: hierarchical pooling across cells with variational inference. HAVI-Methyl differs in that it operates on cfDNA fragmentomic features rather than direct methylation calls, runs amortized inference rather than per-instance variational fitting, and produces calibrated posterior intervals through a conformal wrapper. Among methylation-based deconvolution tools, CelFiE \citep{caggiano2021celfie} performs probabilistic cell-type deconvolution with EM, CelFEER \citep{keukeleire2023celfeer} extends this to read-level deconvolution, and MetDecode \citep{decroos2024metdecode} provides a recent thirteen-entity Bayesian-flavoured deconvolution; these methods consume methylation calls and are downstream of HAVI-Methyl's prediction step. None of the methods consume raw cfDNA WGS, produce calibrated posterior intervals, pool statistical strength hierarchically across samples, and jointly infer tissue-of-origin in a single end-to-end framework; HAVI-Methyl provides all four.

\subsection{cfDNA tumor fraction and copy-number}
A complementary line of work estimates tumor fraction and copy-number directly from cfDNA WGS without going through methylation. ichorCNA \citep{adalsteinsson2017ichor} uses an HMM on read-depth ratios to call copy-number alterations and tumor fraction; PureCN \citep{riester2016purecn} extends this with somatic-variant joint inference. The HAVI-Methyl framework is complementary: ichorCNA-derived tumor fraction can be used as a covariate in the per-sample shift $\delta_s$, and CNV-corrected coverage can be supplied to the Negative-Binomial coverage likelihood. Joint inference of methylation and copy-number is a natural extension and is described in Section~\ref{sec:discussion}.

% ============================================================
